%==============================================================================
\chapter{The Configuration tool}\label{ch:config}
%==============================================================================

The Configuration tool is a helper writing nothing that the matrix pane cannot express by hand, producing a base configuration of frames, gains, delays and crossovers (Chapter~\ref{ch:matrix}). It saves typing (a 5.1 layout is about a dozen crosspoints, a 16-speaker third-order ambisonic decoder up to 256).

For standard layouts the tool writes a whole matrix for you. Choose a layout (\textbf{2.0, 2.1, 4.0, 4.1, 5.1, 7.1}), assign the input and output channel of each speaker, set the gains and the bass-management \gterm{crossover}{crossover} frequency, then apply to a configuration (A..F). The mains are high-passed and their sum is sent, low-passed, to the subwoofer output.

The ambisonics case need a special attention an is exposed in the rest of this chapter.

\section{Ambisonics decoding}\label{sec:ambisonics}
The config tool also generates periphonic (3D) ambisonics decoders for orders 1-3 (Appendix~\ref{ch:ambiprimer} is a primer on the underlying theory). It follows the AmbiX convention (channels in ACN order with SN3D normalisation) and takes the B-format signal on the first $(\mathrm{order}+1)^2$ inputs (4, 9 or 16 channels). Instead of pairing one input with one speaker, each row now describes one loudspeaker: its direction (azimuth and elevation), its output channel, and its distance from the centre. Only the direction takes part in the decode. The distance is not an ambisonic quantity at all, and never reaches the decoder matrix; it drives the optional radius compensation described below, which is a per-output delay and trim for speakers that are not equidistant. The "Speakers" control chooses the number of outputs: at the canonical count the tool ships a hand-tuned periphonic rig (an 8-speaker cube, a 12-speaker icosahedron, a 16-speaker dome); at any other count it lays the speakers out on a near-uniform spherical-Fibonacci lattice. Either way you then edit the angles and radii to match the real installation (the minimum is $(\mathrm{order}+1)^2$ speakers, the number of harmonics).

Clicking "Apply" builds a max-rE sampling decoder: for each speaker $s$ in direction $\Omega_s$ and each harmonic $j$ the decode gain is
\begin{equation}
  D[s][j] \;\propto\; a_{n(j)}\,Y_j^{\mathrm{SN3D}}(\Omega_s),
\end{equation}
where $Y_j^{\mathrm{SN3D}}$ is the real SN3D spherical harmonic of ACN index $j$, $a_n$ the per-order max-rE gains (from the largest root of the Legendre polynomial $P_{\mathrm{order}+1}$, improving localisation), and the matrix is scaled so an on-axis plane wave decodes to about unity. Each non-negligible $D[s][j]$ is written as a harmonic\,$\rightarrow$\,speaker frame, its sign becoming a \gterm{polarity}{polarity} flip. With bass management enabled, the omnidirectional $W$ component (ACN~0) is low-passed to the subwoofer while the speaker feeds are high-passed, which, summed at each output, is exactly the high-passed decoded speaker because the filters are linear. The matrix grows to expose the B-format inputs and the speaker (plus sub) outputs. Refine gains, delays and per-output \gterm{fir}{FIR} corrections in the matrix afterwards.

\paragraph{Importing an IEM AllRADecoder.} As an alternative to the built-in sampling decoder, "Load IEM decoder\ldots" reads a \texttt{.json} configuration exported by the IEM \href{https://plugins.iem.at/#tab-AllRADecoder}{AllRADecoder} (part of the IEM Plug-in Suite). AllRAD (all-round ambisonic decoding) handles irregular and hemispherical rigs better than a plain sampling decoder. The file's decode \texttt{Matrix} replaces the computed one; SuperMoTo converts it to its SN3D\,/\,ACN feed on the fly:
\begin{itemize}[nosep]
  \item if \texttt{ExpectedInputNormalization} is \texttt{n3d}, each coefficient of ACN degree $n$ is multiplied by $\sqrt{2n+1}$ (N3D\,$\to$\,SN3D);
  \item if \texttt{Weights} is \texttt{maxrE} and \texttt{WeightsAlreadyApplied} is false, the per-degree max-rE gain is folded in;
  \item \texttt{Routing} maps each matrix row to a hardware output channel; \texttt{IsImaginary} loudspeakers (used only to regularise the decoder design) carry no row and are skipped;
  \item the loudspeaker layout supplies each speaker's azimuth / elevation, its \textbf{Radius} (used by the radius compensation below) and its per-speaker \textbf{Gain} (folded into the frame gain as a trim).
\end{itemize}
A gain matrix cannot encode time, so the imported decoder, exactly like the built-in one, carries no delay. The radius delay below is the only time alignment, and it is optional.

\paragraph{Radius compensation.} Ambisonic decoding assumes the speakers are equidistant from the centre. When they are not, the tool eventually corrects for each radius $r_s$ relative to the farthest speaker $r_{\max}$, writing a per-speaker output "delay" and "trim":
\begin{equation}
  \Delta t_s = \frac{r_{\max}-r_s}{c}, \qquad
  \Delta g_s = 20\log_{10}\frac{r_s}{r_{\max}} \;\;(\le 0),
\end{equation}
with $c=\SI{343}{\metre\per\second}$. Closer speakers are thus delayed (so all wavefronts reach the centre a the same time) and attenuated, the reference being the farthest speaker. No channel is ever boosted or asked for a negative delay. The "Write radius gain" and "Write radius delay" are independently optionnal, and both apply equally to the built-in and the imported (IEM) decoders. Turning one off leaves that output property untouched, so an alignment or level already set (in the group analysis, or in the single-speaker analysis pane) survives an "Apply". The radius delay is independent of, and rides on top of, the inter-output FIR \gterm{latencycomp}{latency compensation} (Section~\ref{sec:latencycomp}), which only equalises the FIR processing latency. The final arrival of a speaker is $\Delta t_s + l_{\max}$, so the radius alignment survives untouched.

\paragraph{Measured vs.\ geometric alignment (do not double-count).} The geometric radius delay above and a measured alignment are two routes to the same delay, not two delays to add. Both write the same per-output delay, and each simply overwrites the other, so the question is what method you want to determine that delay (estimation based on speaker positions or measurement). There are two ways to measure it. First is the group analysis (Section~\ref{sec:groupalign}), which aligns a whole rig on its most-distant driver and is the natural choice for an ambisonic array. The other way is for a single speaker, the subwoofer time-alignment (Section~\ref{sec:timealign}). Both reference the farthest driver exactly as the radius delay does, and both measure what the radius can only estimate. When you run the per-output sub analysis for a main speaker $k$ at distance $d_k$ from the mic centroid, the main set has $d_k/c$ removed per file and the sub is anchored on that same delay, so the reported "Mains delay" is $\approx (d_{\mathrm{sub}}-d_k)/c$ plus the intrinsic \gterm{groupdelay}{group delay}. Setting that on the output gives an acoustic-plus-electronic arrival of
\begin{equation}
  \underbrace{d_k/c}_{\text{acoustic}} + \underbrace{(d_{\mathrm{sub}}-d_k)/c}_{\text{electronic}}
  = d_{\mathrm{sub}}/c,
\end{equation}
independent of $d_k$: every measured main lands at the same time as the sub and as the other mains. So a per-speaker measured alignment already contains the path-length (radius) term, and more accurately than an entered radius, since it uses the true acoustic path, references the same centroid as the average, and folds in the speakers' intrinsic group delay. Therefore:
\begin{itemize}[nosep]
  \item \textbf{Measured route} (best): align the rig in the group analysis, which writes each output's delay for you. Alternatively, you can do one speaker at a time, by running the sub analysis per main and set each output delay to the reported value. To (re)build the decode afterwards without clobbering those  delays, turn "Write radius delay" off before "Apply" so that the tool then writes the decode, the crossover and the radius trim (unless "Write radius gain" is also off) but leaves each output's delay untouched. (Otherwise set the radii equal so the geometric delay is zero.) Applying the decode with the toggle left on replaces a measured alignment with the geometric estimate, which is the one real way to lose work here.
  \item \textbf{Geometric route} (quick): enter the radii and let the tool write the delay and trim. Use this when there is no sub, or you have not measured each speaker.
\end{itemize}
Level matching is a separate question from alignment which the correction FIR does not answer, beacause each speaker's correction is normalised to that speaker's own mid-band, so inter-speaker amplitude differences survive it untouched. The group analysis does answer it by deriving a suggested level-matching trim per speaker from the corrected mid-band level (Section~\ref{sec:grouplevel}). Like the radius trim it is referenced to the quietest speaker, so it only ever attenuates. It is informational only, though: "Apply \& export" never writes it, so you enter it yourself in that output's \textbf{Trim} in the matrix. Failing that, the radius trim $\Delta g_s$ (or a calibrated-SPL measurement) still handles it.

% TODO(screenshot): the Configuration tool dialog with a 5.1 layout selected,
% per-speaker input/output assignment, gains, and the crossover control.
\screenshot{config-tool.png}{The Configuration tool: a standard layout (here 5.1)
with per-speaker input/output assignment, gains and the bass-management
crossover frequency, ready to Apply into one of the A--F configurations.}
